%!TEX root = ../template.tex
%%%%%%%%%%%%%%%%%%%%%%%%%%%%%%%%%%%%%%%%%%%%%%%%%%%%%%%%%%%%%%%%%%%%
%% abstrac-en.tex
%% ISEL thesis document file
%%
%% Abstract in English
%%%%%%%%%%%%%%%%%%%%%%%%%%%%%%%%%%%%%%%%%%%%%%%%%%%%%%%%%%%%%%%%%%%%
This dissertation is situated in the areas of Computer Engineering and Health, focusing on the development of a recommendation system based on Knowledge Graphs to support clinical research and drug repositioning. The large volume of biomedical data currently available is often distributed across multiple sources, represented in heterogeneous formats, and structured at different levels of detail, which makes its effective integration and analysis challenging. In this context, this dissertation presents a system capable of predicting potential associations between drugs and diseases through link prediction in multi-relational graphs, with the aim of identifying drug repositioning opportunities. Thus, this research investigates how the distinction between different types of entities and relations, explored through Relational Graph Convolutional Networks (R-GCN), contributes to more relevant and interpretable recommendations.

The developed solution is based on the construction of a Knowledge Graph from the integration of regulatory data obtained from DailyMed, Orange Book, and Purple Book. For this purpose, a pipeline for the acquisition, conversion, extraction, and normalization of biomedical entities is used, relying on the MedJsonify framework as a supporting tool for information extraction and mapping to reference ontologies. The resulting graph represents entities such as drugs, active substances, diseases, and genes, as well as relations of composition, therapeutic indication, contraindication, and disease-gene associations, allowing originally dispersed information to be structured in a format suitable for the application of graph-based machine learning models.

On top of this representation, an approach based on Relational Graph Convolutional Networks (R-GCN) is applied, which is suitable for modelling multi-relational graphs by explicitly considering the different types of relations between entities. The effectiveness of the system is quantitatively evaluated using Mean Reciprocal Rank (MRR) and Hits@k metrics, allowing the performance of the model in predicting new links to be analysed. Additionally, a qualitative validation is performed by comparing the generated recommendations with evidence found in scientific literature and clinical trial records. The results suggest that learning over biomedical relational structures contributes to the identification of interpretable and biologically plausible recommendations, supporting the generation of hypotheses for drug repositioning.


\begin{keywords}
Recommendation System, Knowledge Graphs, Relational Graph Convolutional Networks (R-GCN), Link Prediction, Drug Repositioning.
\end{keywords} 
